\documentclass[12pt,a4paper]{article}

\usepackage[utf8]{inputenc}
\usepackage[T1]{fontenc}
\usepackage{geometry}
\usepackage{booktabs}
\usepackage{array}
\usepackage{longtable}
\usepackage{listings}
\usepackage{xcolor}
\usepackage{hyperref}
\usepackage{float}
\usepackage{amsmath}

\geometry{margin=2.5cm}
\hypersetup{
    colorlinks=true,
    linkcolor=black,
    urlcolor=blue,
    citecolor=black
}

\lstdefinestyle{console}{
    basicstyle=\ttfamily\small,
    breaklines=true,
    frame=single,
    backgroundcolor=\color{gray!8},
    columns=fullflexible
}

\lstdefinestyle{cpp}{
    language=C++,
    basicstyle=\ttfamily\small,
    keywordstyle=\color{blue!70!black},
    commentstyle=\color{green!45!black},
    stringstyle=\color{red!60!black},
    breaklines=true,
    frame=single,
    backgroundcolor=\color{gray!8},
    columns=fullflexible
}

\title{\textbf{Memoria de la pr\'actica}\\[0.3cm]
Analizador distribuido de v\'ideo con MPI y GPU}
\author{Asignatura: Arquitecturas Paralelas}
\date{\today}

\begin{document}

\maketitle
\thispagestyle{empty}

\begin{abstract}
Este documento describe la arquitectura, funcionamiento e implementaci\'on del proyecto
\textit{MPI Video Analyzer}, una aplicaci\'on C++ para analizar v\'ideos y detectar personas
por frame mediante modelos YOLO11 exportados a ONNX. El inter\'es principal de la pr\'actica
est\'a en el uso combinado de paralelismo distribuido con MPI y aceleraci\'on de c\'omputo en
GPU mediante CUDA u OpenCL. La aplicaci\'on divide el v\'ideo en lotes de frames, reparte
dichos lotes entre procesos worker, ejecuta el preprocesado y la inferencia en paralelo, y
agrega los resultados en el proceso maestro para producir ficheros CSV, JSON y opcionalmente
un v\'ideo anotado.
\end{abstract}

\newpage
\tableofcontents
\newpage

\section{Introducci\'on}

El objetivo del proyecto es procesar un v\'ideo de entrada para detectar y contar personas en
cada frame. La aplicaci\'on est\'a pensada para escenarios de vigilancia o rescate, donde el
tiempo de procesamiento puede reducirse distribuyendo el trabajo entre varios procesos y,
cuando existen dispositivos disponibles, acelerando las fases m\'as costosas en GPU.

Desde el punto de vista de Arquitecturas Paralelas, el programa combina dos niveles de
paralelismo:

\begin{itemize}
    \item \textbf{Paralelismo de grano grueso con MPI}: el v\'ideo se divide en lotes de frames
    y cada proceso worker procesa lotes independientes.
    \item \textbf{Paralelismo masivo en GPU}: cada worker puede usar CUDA u OpenCL para
    acelerar el preprocesado de imagen y la inferencia del modelo neuronal.
\end{itemize}

El resultado no es solo una aplicaci\'on funcional de visi\'on por computador, sino tambi\'en
una plataforma para estudiar reparto de carga, escalabilidad, sobrecostes de comunicaci\'on,
tama\~no de lote y heterogeneidad entre dispositivos.

\section{Resumen del proyecto}

El ejecutable principal se llama \texttt{rescue\_video\_analyzer}. Se invoca con un v\'ideo de
entrada, un modelo YOLO11 en formato ONNX y varios par\'ametros de ejecuci\'on. Su salida
principal se escribe en el directorio configurado mediante \texttt{--output-dir}.

\begin{itemize}
    \item \texttt{frame\_results.csv}: m\'etricas por frame, incluyendo rank que lo proces\'o,
    GPU usada, conteo de personas y tiempos por etapa.
    \item \texttt{summary.json}: resumen global con metadatos del v\'ideo, configuraci\'on MPI,
    estad\'isticas de detecci\'on, balanceo de carga y tiempos totales.
    \item \texttt{annotated\_people.mp4} o \texttt{annotated\_people.avi}: v\'ideo anotado con
    cajas de detecci\'on y conteo visual, salvo que se use \texttt{--no-annotated-video}.
\end{itemize}

\section{Tecnolog\'ias utilizadas}

La pr\'actica integra varias tecnolog\'ias de computaci\'on paralela y visi\'on artificial.

\begin{longtable}{p{0.25\textwidth}p{0.65\textwidth}}
\toprule
\textbf{Tecnolog\'ia} & \textbf{Uso en el proyecto} \\
\midrule
C++17 & Lenguaje principal de la aplicaci\'on. \\
OpenMPI & Ejecuci\'on distribuida mediante procesos MPI. \\
CUDA & Backend principal para preprocesado GPU e inferencia con ONNX Runtime CUDA. \\
OpenCL & Backend alternativo mediante OpenCV \texttt{UMat} y OpenCV DNN. \\
OpenCV & Lectura/escritura de v\'ideo, transformaciones de imagen, NMS y generaci\'on del v\'ideo anotado. \\
ONNX Runtime & Ejecuci\'on del modelo YOLO11 en el backend CUDA. \\
YOLO11 ONNX & Modelo de detecci\'on de objetos; se filtra la clase persona. \\
CMake & Sistema de compilaci\'on del proyecto. \\
\bottomrule
\end{longtable}

\section{Organizaci\'on del c\'odigo}

La estructura principal del repositorio es la siguiente:

\begin{lstlisting}[style=console]
.
|-- CMakeLists.txt
|-- README.md
|-- dataset/                 # videos de prueba
|-- include/rescue/          # cabeceras publicas
|-- models/                  # modelos YOLO11 en ONNX
|-- scripts/
|   |-- benchmark_parallel.sh
|   `-- launch_cluster.py
`-- src/                     # implementacion C++, CUDA y OpenCL
\end{lstlisting}

Los m\'odulos m\'as importantes son:

\begin{itemize}
    \item \texttt{src/main.cpp}: parseo de argumentos, preparaci\'on del entorno y arranque MPI.
    \item \texttt{src/mpi\_runtime.cpp}: l\'ogica maestro-worker, planificadores y comunicaci\'on MPI.
    \item \texttt{src/gpu\_preprocess.cu}: kernels CUDA de preprocesado.
    \item \texttt{src/gpu\_preprocess\_opencl.cpp}: preprocesado alternativo con OpenCV/OpenCL.
    \item \texttt{src/yolo\_detector\_cuda.cpp}: inferencia YOLO11 con ONNX Runtime CUDA.
    \item \texttt{src/yolo\_detector\_opencl.cpp}: inferencia YOLO11 con OpenCV DNN y OpenCL.
    \item \texttt{src/output\_writer.cpp}: generaci\'on de CSV, JSON y v\'ideo anotado.
    \item \texttt{src/tracker.cpp}: tracker sencillo usado solo para visualizaci\'on.
\end{itemize}

\section{Arquitectura paralela}

\subsection{Modelo SPMD}

La aplicaci\'on sigue un modelo SPMD (\textit{Single Program, Multiple Data}): todos los
procesos ejecutan el mismo binario, pero su comportamiento depende del rank MPI.

\begin{itemize}
    \item El \textbf{rank 0} act\'ua como proceso maestro.
    \item Los \textbf{ranks 1..N} act\'uan como workers.
\end{itemize}

El maestro no ejecuta inferencia. Su trabajo principal es leer frames, construir lotes,
enviarlos a los workers, recibir resultados y escribir las salidas. Los workers reciben
lotes, ejecutan el preprocesado y la detecci\'on, y devuelven un paquete de resultados por
frame.

\subsection{Descomposici\'on del problema}

El problema se descompone por frames. Cada frame puede procesarse de forma independiente
para obtener detecciones de personas, por lo que existe paralelismo de datos natural. En vez
de enviar frames individuales, el programa agrupa frames en lotes de tama\~no
\texttt{batch\_size}. Esto reduce el n\'umero de mensajes MPI y permite amortizar parte del
coste de comunicaci\'on.

Si el v\'ideo tiene $F$ frames y el tama\~no de lote es $B$, el n\'umero aproximado de tareas es:

\[
    T = \left\lceil \frac{F}{B} \right\rceil
\]

Cada tarea contiene un \texttt{BatchHeader} y un bloque de bytes con los frames codificados
como matrices \texttt{CV\_8UC3}.

\subsection{Topolog\'ia maestro-worker}

La comunicaci\'on MPI utiliza etiquetas separadas para cabeceras y cargas:

\begin{table}[H]
\centering
\begin{tabular}{lll}
\toprule
\textbf{Etiqueta} & \textbf{Sentido} & \textbf{Contenido} \\
\midrule
\texttt{kTagBatchHeader} & Maestro $\rightarrow$ worker & Cabecera del lote. \\
\texttt{kTagBatchPayload} & Maestro $\rightarrow$ worker & Bytes de los frames. \\
\texttt{kTagResultHeader} & Worker $\rightarrow$ maestro & Cabecera del resultado. \\
\texttt{kTagResultPayload} & Worker $\rightarrow$ maestro & Array de \texttt{FramePacket}. \\
\bottomrule
\end{tabular}
\caption{Mensajes MPI principales.}
\end{table}

Las estructuras simples se env\'ian como tipos MPI contiguos de bytes, ya que
\texttt{BatchHeader}, \texttt{ResultHeader}, \texttt{FrameResult} y \texttt{FramePacket} son
trivialmente copiables. Esta decisi\'on simplifica el transporte de datos estructurados, aunque
presupone que los procesos ejecutan una versi\'on binaria compatible de la aplicaci\'on.

\subsection{Planificadores}

El proyecto implementa tres modos de scheduling:

\begin{itemize}
    \item \textbf{\texttt{dynamic}}: el maestro entrega un nuevo lote al worker que acaba de
    terminar. Es el modo recomendado en entornos heterog\'eneos porque se adapta al ritmo real
    de cada GPU.
    \item \textbf{\texttt{static-contiguous}}: cada worker recibe un bloque contiguo de frames.
    Reduce decisiones en tiempo de ejecuci\'on, pero puede producir desequilibrio si unos
    frames son m\'as costosos que otros o si los dispositivos no tienen el mismo rendimiento.
    \item \textbf{\texttt{static-round-robin}}: los lotes se asignan por turnos a los workers.
    Reparte mejor que el bloque contiguo cuando el coste cambia a lo largo del v\'ideo, aunque
    no reacciona al rendimiento real de cada worker.
\end{itemize}

El scheduling din\'amico es especialmente relevante en Arquitecturas Paralelas porque reduce
el tiempo ocioso de los workers. Cuando un worker termina, el maestro recibe su resultado con
\texttt{MPI\_Probe(MPI\_ANY\_SOURCE, ...)} y le asigna trabajo nuevo si quedan lotes pendientes.

\section{Flujo de ejecuci\'on}

\subsection{Arranque}

El programa valida los argumentos de l\'inea de comandos y comprueba que existan tanto el
v\'ideo de entrada como el modelo ONNX. Tambi\'en prepara rutas de bibliotecas din\'amicas para
ONNX Runtime y CUDA, especialmente en entornos WSL o virtualenv.

Si el binario se ejecuta sin \texttt{mpirun}, puede relanzarse autom\'aticamente con un mundo
MPI por defecto. En CUDA, el esquema natural es:

\[
    \texttt{world\_size} = 1 + \texttt{numero\_de\_GPUs}
\]

donde el proceso adicional es el maestro.

\subsection{Trabajo del maestro}

El maestro realiza los siguientes pasos:

\begin{enumerate}
    \item Abre el v\'ideo con OpenCV y obtiene metadatos: resoluci\'on, FPS, n\'umero de frames y duraci\'on.
    \item Construye un plan de lotes seg\'un \texttt{batch-size}.
    \item Env\'ia un primer lote a cada worker disponible.
    \item Espera resultados de cualquier worker.
    \item Inserta los \texttt{FramePacket} recibidos en un vector global.
    \item Env\'ia un nuevo lote al worker que ha quedado libre, o una se\~nal de parada si ya no hay trabajo.
    \item Ordena los resultados por \texttt{frame\_index}.
    \item Escribe las salidas CSV, JSON y v\'ideo anotado.
\end{enumerate}

Este dise\~no centraliza la entrada/salida de v\'ideo en el maestro y evita que varios procesos
lean simult\'aneamente el mismo fichero. A cambio, el maestro se convierte en punto de
coordinaci\'on y en posible cuello de botella si el n\'umero de workers crece mucho.

\subsection{Trabajo de cada worker}

Cada worker ejecuta un bucle hasta recibir una cabecera de lote con \texttt{frame\_count = 0}.
Para cada lote:

\begin{enumerate}
    \item Recibe la cabecera y el payload de frames.
    \item Materializa vistas \texttt{cv::Mat} sobre el bloque de bytes recibido.
    \item Realiza una fase de calentamiento del backend si todav\'ia no se ha hecho.
    \item Redimensiona cada frame a la resoluci\'on de procesamiento si es necesario.
    \item Ejecuta el preprocesado GPU: \textit{letterbox}, normalizaci\'on, conversi\'on a tensor y m\'etricas auxiliares.
    \item Ejecuta la inferencia YOLO11.
    \item Filtra detecciones de clase persona y aplica NMS cuando corresponde.
    \item Construye un \texttt{FramePacket} por frame y lo devuelve al maestro.
\end{enumerate}

\section{Backends de GPU}

\subsection{Ruta CUDA}

La ruta CUDA es la ruta principal del proyecto. Se compila si \texttt{RESCUE\_ENABLE\_CUDA=ON}
en CMake. En esta ruta, cada worker activa una GPU mediante \texttt{cudaSetDevice}. Si no se
fuerza un dispositivo con \texttt{--gpu-device}, el programa usa el rank local del worker para
repartir workers entre GPUs visibles:

\[
    \texttt{device\_id} =
    \texttt{local\_worker\_rank} \bmod \texttt{visible\_device\_count}
\]

El preprocesado CUDA se implementa en \texttt{src/gpu\_preprocess.cu}. El kernel principal:

\begin{itemize}
    \item Lee la imagen BGR de entrada.
    \item Calcula el reescalado con \textit{letterbox}, manteniendo la relaci\'on de aspecto.
    \item Aplica interpolaci\'on bilineal.
    \item Convierte BGR a RGB.
    \item Normaliza los valores a $[0,1]$.
    \item Escribe el tensor en formato \texttt{NCHW}, esperado por YOLO.
    \item Genera una imagen en gris auxiliar para calcular intensidad media y densidad de bordes.
\end{itemize}

Un segundo kernel calcula estad\'isticas con operaciones at\'omicas: suma de intensidades y
conteo de p\'ixeles cuyo gradiente supera el umbral \texttt{edge-threshold}. La inferencia se
realiza con ONNX Runtime y el \textit{CUDA Execution Provider}. El tensor de entrada se crea
directamente sobre memoria CUDA, evitando copiarlo a CPU antes de ejecutar el modelo.

\subsection{Ruta OpenCL}

La ruta OpenCL es una alternativa basada en OpenCV. Usa \texttt{cv::UMat} para permitir que
OpenCV ejecute operaciones sobre OpenCL cuando est\'a disponible. El preprocesado incluye
redimensionado, \textit{letterbox}, creaci\'on del blob y c\'alculo de m\'etricas de imagen con
Sobel. La inferencia se realiza con \texttt{cv::dnn::readNetFromONNX} y target
\texttt{DNN\_TARGET\_OPENCL}. Si OpenCL no est\'a disponible, OpenCV puede caer a CPU.

\subsection{Backend autom\'atico}

Con \texttt{--gpu-backend auto}, el worker prefiere CUDA si hay dispositivos visibles y el
proyecto fue compilado con soporte CUDA. En caso contrario, utiliza la ruta OpenCL/CPU. Este
comportamiento facilita ejecutar el mismo binario en equipos con distintas capacidades.

\section{Modelo de detecci\'on YOLO11}

Los modelos incluidos en \texttt{models/} est\'an exportados a ONNX. El modelo por defecto es:

\begin{lstlisting}[style=console]
models/yolo11s_1088.onnx
\end{lstlisting}

El par\'ametro \texttt{--resize-width} debe coincidir con el tama\~no de entrada usado al
exportar el modelo. Por ejemplo, \texttt{yolo11s\_1088.onnx} se utiliza con
\texttt{--resize-width 1088}. Tras la inferencia, el programa decodifica la salida de YOLO,
filtra la clase persona (\texttt{class\_id = 0}) y descarta detecciones con confianza inferior
a \texttt{--score-threshold}. En salidas no end-to-end se aplica \textit{Non-Maximum
Suppression} con \texttt{--nms-threshold} y \texttt{--top-k}.

\section{Estructuras de datos}

\begin{longtable}{p{0.25\textwidth}p{0.65\textwidth}}
\toprule
\textbf{Estructura} & \textbf{Funci\'on} \\
\midrule
\texttt{AppConfig} & Configuraci\'on global: rutas, tama\~no de lote, backend GPU, modelo, thresholds y scheduler. \\
\texttt{VideoMetadata} & Metadatos del v\'ideo: resoluci\'on, FPS, frames y duraci\'on. \\
\texttt{BatchHeader} & Cabecera enviada por MPI para describir un lote de frames. \\
\texttt{FrameResult} & Resultado num\'erico de un frame: rank, GPU, conteo y tiempos. \\
\texttt{DetectionBox} & Caja de detecci\'on con coordenadas, confianza y clase. \\
\texttt{FramePacket} & Paquete completo devuelto al maestro: \texttt{FrameResult} y hasta 128 detecciones. \\
\bottomrule
\end{longtable}

El l\'imite \texttt{kMaxDetectionsPerFrame = 128} acota el tama\~no del paquete MPI por frame.
Esto simplifica la comunicaci\'on porque todos los \texttt{FramePacket} tienen tama\~no fijo.

\section{M\'etricas y salidas}

La aplicaci\'on registra tiempos por frame:

\begin{itemize}
    \item \texttt{read\_ms}: coste medio asociado a recibir el lote desde MPI.
    \item \texttt{preprocess\_ms}: tiempo de redimensionado/preprocesado.
    \item \texttt{detection\_ms}: tiempo de inferencia y decodificaci\'on.
    \item \texttt{processing\_ms}: suma de recepci\'on, preprocesado e inferencia por frame.
\end{itemize}

En \texttt{summary.json} tambi\'en se escriben m\'etricas globales:

\begin{itemize}
    \item \texttt{distributed\_processing\_ms}: tiempo desde que el maestro empieza a repartir trabajo hasta que recibe todos los resultados.
    \item \texttt{output\_write\_ms}: tiempo de escritura de salidas.
    \item \texttt{total\_wall\_ms}: suma de procesamiento distribuido y escritura.
    \item \texttt{load\_balance}: frames procesados por worker, m\'inimo, m\'aximo y ratio m\'aximo/media.
    \item \texttt{worker\_stats}: estad\'isticas desglosadas por rank.
\end{itemize}

Estas m\'etricas permiten analizar si el programa escala, si hay workers ociosos, si el cuello
de botella est\'a en la inferencia, en el preprocesado, en la comunicaci\'on MPI o en la
escritura de resultados.

\section{Compilaci\'on}

La compilaci\'on normal usa CMake:

\begin{lstlisting}[style=console]
cmake -S . -B build
cmake --build build -j
\end{lstlisting}

Para compilar sin CUDA y usar solo OpenCL/CPU:

\begin{lstlisting}[style=console]
cmake -S . -B build -DRESCUE_ENABLE_CUDA=OFF
cmake --build build -j
\end{lstlisting}

La arquitectura CUDA por defecto en el proyecto es \texttt{89}, adecuada para la GPU usada
durante el desarrollo. Puede cambiarse con:

\begin{lstlisting}[style=console]
cmake -S . -B build -DCMAKE_CUDA_ARCHITECTURES=86
cmake --build build -j
\end{lstlisting}

\section{Ejecuci\'on}

La ejecuci\'on recomendada en una m\'aquina con una GPU NVIDIA es:

\begin{lstlisting}[style=console]
mpirun -np 2 ./build/rescue_video_analyzer \
  --input dataset/videoset2.mp4 \
  --output-dir output \
  --batch-size 8 \
  --scheduler dynamic
\end{lstlisting}

Para medir rendimiento sin incluir el coste de generar v\'ideo anotado:

\begin{lstlisting}[style=console]
mpirun -np 2 ./build/rescue_video_analyzer \
  --input dataset/videoset2.mp4 \
  --output-dir output_bench \
  --batch-size 8 \
  --scheduler dynamic \
  --no-annotated-video
\end{lstlisting}

El modo cluster puede lanzarse con el script:

\begin{lstlisting}[style=console]
./scripts/launch_cluster.py \
  --ips localhost,192.168.1.55 \
  --args "--input dataset/videoset2.mp4 --output-dir output_cluster --batch-size 8 --scheduler dynamic"
\end{lstlisting}

Este script sondea GPUs NVIDIA y OpenCL en los nodos, genera un \texttt{hosts\_cluster.txt}
y lanza \texttt{mpiexec}. En un cluster real, las rutas de \texttt{dataset/}, \texttt{models/}
y del ejecutable deben estar disponibles en todos los nodos o en un sistema de ficheros
compartido.

\section{Benchmarking}

El script \texttt{scripts/benchmark\_parallel.sh} automatiza pruebas variando el tama\~no del
mundo MPI y el scheduler:

\begin{lstlisting}[style=console]
MPI_WORLD_SIZES="2 3 4" \
SCHEDULERS="dynamic static-contiguous static-round-robin" \
./scripts/benchmark_parallel.sh dataset/videoset.mp4 parallel_benchmarks
\end{lstlisting}

El script escribe una tabla TSV con:

\begin{itemize}
    \item \texttt{world\_size}
    \item \texttt{scheduler}
    \item \texttt{total\_wall\_ms}
    \item \texttt{distributed\_processing\_ms}
    \item \texttt{average\_detection\_ms}
    \item \texttt{average\_processing\_ms}
    \item \texttt{average\_people\_per\_frame}
    \item directorio de salida
\end{itemize}

Una prueba local recogida en el README, usando \texttt{dataset/videoset2.mp4},
\texttt{mpirun -np 2}, \texttt{--batch-size 8} y \texttt{--no-annotated-video}, obtuvo:

\begin{table}[H]
\centering
\begin{tabular}{lr}
\toprule
\textbf{M\'etrica} & \textbf{Valor} \\
\midrule
Personas medias por frame & 29.863 \\
Preprocesado medio & 3.841 ms \\
Detecci\'on media & 16.158 ms \\
Procesamiento medio & 20.090 ms \\
Tiempo total & 6399.037 ms \\
\bottomrule
\end{tabular}
\caption{Ejemplo de resultado local documentado.}
\end{table}

Es importante interpretar estos datos con cuidado. En una m\'aquina con una sola GPU, ejecutar
\texttt{np=3} o \texttt{np=4} no implica escalado multi-GPU real: varios workers competir\'an
por el mismo dispositivo o caer\'an a rutas alternativas. Para medir la ruta CUDA principal en
una sola GPU, la configuraci\'on m\'as limpia es \texttt{np=2}: un maestro y un worker.

\section{An\'alisis desde Arquitecturas Paralelas}

\subsection{Granularidad}

La unidad de trabajo es el lote de frames. Un lote demasiado peque\~no incrementa el n\'umero de
mensajes MPI y el coste de scheduling. Un lote demasiado grande reduce la capacidad del
scheduler din\'amico para balancear carga, porque un worker lento puede quedarse con una tarea
grande durante demasiado tiempo. El par\'ametro \texttt{--batch-size} permite estudiar ese
compromiso.

\subsection{Comunicaci\'on frente a c\'omputo}

Cada lote requiere enviar al worker los bytes completos de los frames. Para un frame
\texttt{CV\_8UC3}, el tama\~no aproximado es:

\[
    S = ancho \cdot alto \cdot 3 \text{ bytes}
\]

Por tanto, el coste de comunicaci\'on crece con la resoluci\'on de procesamiento y el tama\~no
del lote. Sin embargo, el coste dominante suele estar en inferencia YOLO, especialmente con
modelos grandes o resoluciones de entrada altas. El dise\~no busca que la comunicaci\'on quede
amortizada por el coste computacional de cada lote.

\subsection{Balanceo de carga}

El modo din\'amico mejora el balanceo cuando:

\begin{itemize}
    \item Hay GPUs con distinto rendimiento.
    \item Alg\'un nodo cae a OpenCL/CPU.
    \item El coste de detecci\'on var\'ia por frame.
    \item Hay contenci\'on por compartir dispositivo.
\end{itemize}

El fichero \texttt{summary.json} incluye el ratio:

\[
    \texttt{max\_to\_mean\_ratio} =
    \frac{\text{frames del worker con m\'as carga}}
         {\text{media de frames por worker activo}}
\]

Un valor cercano a 1 indica buen balanceo en n\'umero de frames. Aun as\'i, el n\'umero de
frames no siempre equivale a tiempo: dos workers pueden procesar el mismo n\'umero de frames
con tiempos distintos si sus GPUs no son equivalentes.

\subsection{Heterogeneidad}

El proyecto contempla heterogeneidad de hardware. El backend autom\'atico decide entre CUDA y
OpenCL/CPU seg\'un disponibilidad. Adem\'as, el script de cluster asigna un slot por GPU
detectada y un slot adicional para el maestro en el nodo inicial. Esta aproximaci\'on permite
usar varios nodos con distinta configuraci\'on, aunque la eficiencia final depende de que el
scheduler mantenga ocupados los dispositivos r\'apidos sin quedar bloqueado por los lentos.

\subsection{Cuellos de botella}

Los posibles cuellos de botella son:

\begin{itemize}
    \item \textbf{Maestro}: lee el v\'ideo, serializa lotes y recibe resultados. Con muchos
    workers puede limitar la escalabilidad.
    \item \textbf{Comunicaci\'on MPI}: el env\'io de frames completos puede ser costoso en red.
    \item \textbf{GPU}: la inferencia YOLO es la fase m\'as pesada por frame.
    \item \textbf{Memoria}: hay copias CPU-GPU en preprocesado CUDA y almacenamiento temporal
    de frames.
    \item \textbf{Salida}: la generaci\'on de v\'ideo anotado puede distorsionar benchmarks, por
    eso existe \texttt{--no-annotated-video}.
\end{itemize}

\section{Limitaciones}

\begin{itemize}
    \item El maestro centraliza la lectura del v\'ideo, lo que simplifica la coherencia pero
    puede limitar escalabilidad en clusters grandes.
    \item Los frames se env\'ian completos por MPI; no hay compresi\'on ni lectura distribuida
    desde cada nodo.
    \item \texttt{FramePacket} fija un m\'aximo de 128 detecciones por frame para facilitar la
    comunicaci\'on.
    \item El tracking no se usa para modificar el c\'omputo distribuido: solo para el v\'ideo
    anotado y el acumulado visual.
    \item La ruta OpenCL depende de la disponibilidad real de OpenCL en OpenCV; si no existe,
    puede ejecutarse en CPU.
\end{itemize}

\section{Posibles mejoras}

\begin{itemize}
    \item Usar comunicaci\'on no bloqueante con \texttt{MPI\_Isend} y \texttt{MPI\_Irecv} para
    solapar comunicaci\'on y c\'omputo.
    \item Permitir lectura distribuida del v\'ideo si el formato y el almacenamiento compartido
    lo permiten.
    \item Enviar frames ya redimensionados para reducir bytes transmitidos, comparando el coste
    de mover preprocesado al maestro frente a mantenerlo en workers.
    \item Agrupar inferencias por batch real en GPU para aprovechar mejor el throughput del
    modelo.
    \item Implementar afinidad m\'as fina entre ranks, NUMA y GPUs en nodos multi-GPU.
    \item Registrar trazas temporales completas para representar gr\'aficamente ocupaci\'on de
    workers y fases de pipeline.
\end{itemize}

\section{Conclusiones}

El proyecto implementa una arquitectura paralela h\'ibrida adecuada para una pr\'actica de
Arquitecturas Paralelas. A nivel distribuido, usa MPI para repartir lotes de frames entre
workers independientes. A nivel de dispositivo, emplea CUDA u OpenCL para acelerar el
preprocesado y la inferencia de un modelo YOLO11. El dise\~no permite comparar schedulers,
medir balanceo de carga y observar la relaci\'on entre comunicaci\'on, c\'omputo y escritura de
resultados.

La principal fortaleza del sistema es que el problema tiene paralelismo de datos natural: los
frames pueden procesarse de forma independiente y agregarse posteriormente ordenando por
\texttt{frame\_index}. El scheduler din\'amico aprovecha esta independencia para adaptarse a
hardware heterog\'eneo. Como limitaci\'on, el maestro concentra la lectura y distribuci\'on de
frames, por lo que una mejora futura clara ser\'ia estudiar t\'ecnicas de solapamiento,
comunicaci\'on no bloqueante y lectura distribuida.

\appendix
\section{Par\'ametros principales}

\begin{longtable}{p{0.35\textwidth}p{0.55\textwidth}}
\toprule
\textbf{Par\'ametro} & \textbf{Descripci\'on} \\
\midrule
\texttt{--input <video>} & V\'ideo de entrada. Obligatorio. \\
\texttt{--output-dir <dir>} & Directorio de salida. \\
\texttt{--batch-size <N>} & Frames por lote MPI. \\
\texttt{--processing-width <px>} & Ancho de procesamiento; 0 usa resoluci\'on nativa. \\
\texttt{--processing-height <px>} & Alto de procesamiento; 0 usa resoluci\'on nativa. \\
\texttt{--gpu-device auto|<id>} & Selecci\'on de GPU visible. \\
\texttt{--gpu-backend auto|cuda|opencl} & Backend de c\'omputo usado por los workers. \\
\texttt{--resize-width <px>} & Tama\~no cuadrado de entrada YOLO. \\
\texttt{--edge-threshold <N>} & Umbral para densidad de bordes auxiliar. \\
\texttt{--model <path>} & Modelo ONNX. \\
\texttt{--score-threshold <0..1>} & Confianza m\'inima de detecci\'on. \\
\texttt{--nms-threshold <0..1>} & Umbral de NMS. \\
\texttt{--top-k <N>} & M\'aximo de candidatos antes de NMS. \\
\texttt{--scheduler <modo>} & \texttt{dynamic}, \texttt{static-contiguous} o \texttt{static-round-robin}. \\
\texttt{--no-annotated-video} & Desactiva generaci\'on del v\'ideo anotado. \\
\bottomrule
\end{longtable}

\section{Dependencias}

En Ubuntu/WSL, las dependencias principales son:

\begin{lstlisting}[style=console]
sudo apt-get update
sudo apt-get install -y \
  cmake \
  libopencv-dev \
  openmpi-bin \
  libopenmpi-dev \
  nvidia-cuda-toolkit
\end{lstlisting}

\end{document}
